\chapter{Átomos polielectrónicos y símbolos de término}\label{ch:b3:many-electron-atoms}

En 1868, una raya amarilla del espectro del Sol reveló un elemento
desconocido en la Tierra, el helio. Cuando por fin se aisló en el laboratorio, en
1895, su espectro parecía el de \emph{dos} elementos: dos familias de
rayas, atribuidas durante un tiempo a un ``ortohelio'' y a un
``parahelio'', y que nunca intercambiaban luz entre sí. Solo hay
un helio. Las dos familias son los \omterm{def:b3:many-electron-atoms:singlet-triplet}{estados triplete} y singlete del
mismo átomo, separados por el espín del electrón y por una regla más fundamental
que cualquier fuerza: la función de onda de dos electrones debe cambiar de signo cuando
se intercambian. Este capítulo sigue esa regla desde el principio de Pauli hasta
los \omterm{def:b3:many-electron-atoms:term}{símbolos de término} con los que los espectroscopistas etiquetan cada estado de cada átomo
e ion ---incluidos los iones de metales de transición cuyos colores se explican
en el \cref{ch:b3:complex-spectra-magnetism}.

\begin{recall}
El volumen del primer año etiquetó los electrones con cuatro números cuánticos $n$, $l$,
$m_l$, $m_s$, construyó configuraciones fundamentales con el principio de Pauli, la
regla de Klechkowski y la regla de Hund, y contó electrones desapareados. El volumen del segundo
año dio las energías orbitales y las reglas de Slater del apantallamiento.
El \cref{ch:b3:quantum-model-systems} aportó los \omterm{def:b3:quantum-model-systems:operator}{operadores}, los \omterm{def:b3:quantum-model-systems:operator}{valores propios}, el
momento angular del rotor ($\hat L^2$ con \omterm{def:b3:quantum-model-systems:operator}{valores propios} $l(l+1)\hbar^2$) y el
átomo de hidrógeno.
\end{recall}

\section{Espín y espín-orbitales}

El electrón posee un momento angular intrínseco, su espín, de número
cuántico $s = \frac12$: $\hat S^2$ tiene el único \omterm{def:b3:quantum-model-systems:operator}{valor propio} $s(s+1)\hbar^2 =
\frac34\hbar^2$, y $\hat S_z$ los dos \omterm{def:b3:quantum-model-systems:operator}{valores propios} $m_s\hbar = \pm\frac12\hbar$.
Las dos funciones de espín se escriben $\alpha$ ($m_s = +\frac12$) y $\beta$
($m_s = -\frac12$); son ortonormales, $\langle\alpha|\alpha\rangle =
\langle\beta|\beta\rangle = 1$ y $\langle\alpha|\beta\rangle = 0$. El espín no tiene
equivalente clásico; no es la rotación de una pequeña esfera, y no
interviene en absoluto en un hamiltoniano no relativista.

\begin{definition}[Espín-orbital]\label{def:b3:many-electron-atoms:spin-orbital}
Un \emph{espín-orbital}\index{espin-orbital@espín-orbital} es el producto de un orbital espacial
$\phi(\mathbf r)$ por una función de espín: $\phi\alpha$ o $\phi\beta$. Un
espín-orbital contiene la descripción completa de un electrón.
\end{definition}

Un orbital atómico etiquetado $(n, l, m_l)$ en el volumen del primer año da así dos
\omterm{def:b3:many-electron-atoms:spin-orbital}{espín-orbitales}, y por eso una subcapa aloja $2(2l+1)$ electrones. La verdadera
pregunta es por qué un \omterm{def:b3:many-electron-atoms:spin-orbital}{espín-orbital} aloja como máximo uno.

\section{Electrones indistinguibles y determinante de Slater}

\begin{definition}[Partículas indistinguibles, función de onda antisimétrica]\label{def:b3:many-electron-atoms:indistinguishable}
Unas partículas son \emph{indistinguibles}\index{particulas indistinguibles@partículas indistinguibles} si
ninguna medida permite diferenciarlas: todo electrón es idéntico a cualquier
otro. Una función de onda de varios electrones es una \emph{función de onda
antisimétrica}\index{funcion de onda antisimetrica@función de onda antisimétrica} si cambia de signo cuando se
intercambian las coordenadas (de espacio y de espín) de dos electrones cualesquiera:
$\Psi(\dots,i,\dots,j,\dots) = -\Psi(\dots,j,\dots,i,\dots)$.
\end{definition}

La indistinguibilidad por sí sola exige que $|\Psi|^2$ no cambie en un
intercambio, de modo que $\Psi$ puede, como mucho, cambiar de signo. La naturaleza ha elegido: la
función de onda de cualquier sistema de electrones ---de cualquier conjunto de fermiones--- es
antisimétrica. Es un postulado de la mecánica cuántica, confirmado por todos los
espectros atómicos.

\begin{definition}[Determinante de Slater]\label{def:b3:many-electron-atoms:slater}
Para $N$ electrones en $N$ \omterm{def:b3:many-electron-atoms:spin-orbital}{espín-orbitales} distintos $\chi_1, \dots, \chi_N$, el
\emph{determinante de Slater}\index{determinante de Slater} es
\[
\Psi = \frac{1}{\sqrt{N!}}\begin{vmatrix}
\chi_1(1) & \chi_2(1) & \cdots & \chi_N(1)\\
\chi_1(2) & \chi_2(2) & \cdots & \chi_N(2)\\
\vdots & & & \vdots\\
\chi_1(N) & \chi_2(N) & \cdots & \chi_N(N)
\end{vmatrix},
\]
donde la fila $i$ contiene el electrón $i$ en cada \omterm{def:b3:many-electron-atoms:spin-orbital}{espín-orbital}. Se escribe abreviadamente
$|\chi_1\chi_2\dots\chi_N|$.
\end{definition}

\begin{theorem}[El principio de Pauli]\label{thm:b3:many-electron-atoms:pauli}
Un \omterm{def:b3:many-electron-atoms:slater}{determinante de Slater} es antisimétrico, y se anula si dos electrones
ocupan el mismo \omterm{def:b3:many-electron-atoms:spin-orbital}{espín-orbital}. Por tanto, dos electrones de un átomo no pueden tener los
mismos cuatro números cuánticos.
\end{theorem}

\begin{proof}
Intercambiar los electrones $i$ y $j$ intercambia dos filas del determinante,
lo que cambia su signo: la antisimetría se cumple para cada par. Si dos
\omterm{def:b3:many-electron-atoms:spin-orbital}{espín-orbitales} son iguales, dos columnas son iguales y el determinante es nulo:
no existe tal estado. Dos electrones del mismo átomo con los mismos $n$, $l$,
$m_l$ y $m_s$ ocuparían el mismo \omterm{def:b3:many-electron-atoms:spin-orbital}{espín-orbital}.
\end{proof}

\begin{example}[El estado fundamental del helio]\label{ex:b3:many-electron-atoms:helium-ground}
Con los dos electrones en $1s$,
\[
\Psi = \frac{1}{\sqrt2}\begin{vmatrix}1s\alpha(1) & 1s\beta(1)\\ 1s\alpha(2) &
1s\beta(2)\end{vmatrix} = 1s(1)\,1s(2)\cdot\frac{1}{\sqrt2}\big[\alpha(1)\beta(2)
- \beta(1)\alpha(2)\big].
\]
La parte espacial es simétrica y la parte de espín antisimétrica: los dos espines
están apareados, con espín total nulo.
\end{example}

\section{Helio: repulsión coulombiana e intercambio}

Se excita un electrón del helio a $2s$. Con un electrón en cada orbital, pueden construirse cuatro
determinantes a partir de $1s\alpha$, $1s\beta$, $2s\alpha$, $2s\beta$.
Reordenados, se convierten en productos de una función espacial por una función de espín:
\[
\Psi_\pm = \frac{1}{\sqrt2}\big[1s(1)2s(2) \pm 2s(1)1s(2)\big]\times(\text{función
de espín}).
\]
La función espacial antisimétrica $\Psi_-$ debe acompañarse de una función de espín
simétrica, de las que hay tres: $\alpha(1)\alpha(2)$,
$\beta(1)\beta(2)$ y $\frac{1}{\sqrt2}[\alpha(1)\beta(2) + \beta(1)\alpha(2)]$, las
tres componentes $M_S = 1, 0, -1$ de un espín total $S = 1$. La $\Psi_+$
simétrica debe acompañarse de la única función de espín antisimétrica, $S = 0$.

\begin{definition}[Estados singlete y triplete]\label{def:b3:many-electron-atoms:singlet-triplet}
Un estado de espín total $S = 0$ es un \emph{estado singlete}\index{estado singlete};
uno de espín total $S = 1$, cuyas tres componentes $M_S = 1, 0, -1$ tienen la
misma energía en ausencia de campo magnético, es un \emph{estado
triplete}\index{estado triplete}.
\end{definition}

\begin{definition}[Integral de intercambio]\label{def:b3:many-electron-atoms:exchange}
Para dos orbitales $a$ y $b$ y la repulsión electrónica $e^2/4\pi\varepsilon_0
r_{12}$, la \emph{integral de intercambio}\index{integral de intercambio} es
\[
K_{ab} = \iint a(1)b(2)\,\frac{e^2}{4\pi\varepsilon_0r_{12}}\,b(1)a(2)\,\dd\tau_1\dd\tau_2,
\]
la misma integral que la repulsión electrónica clásica $J_{ab}$ (con
$a(1)b(2)$ a ambos lados) salvo que los electrones están intercambiados en un
lado. $K_{ab}$ es positiva y no tiene equivalente clásico.
\end{definition}

\begin{proposition}[Energías del singlete y del triplete]\label{prop:b3:many-electron-atoms:helium}
A primer orden en la repulsión electrónica, los estados $1s2s$ del helio tienen
las energías $E_\pm = E_{1s} + E_{2s} + J \pm K$, el singlete ($+$) por encima del
triplete ($-$) en $2K$.
\end{proposition}

\begin{proof}
La energía no perturbada de ambas $\Psi_\pm$ es $E_{1s} + E_{2s}$ (orbitales
hidrogenoides de carga nuclear 2). La corrección de primer orden es el
\omterm{def:b3:quantum-model-systems:hermitian}{valor esperado} de la repulsión $\hat V_{12}$ en el estado normalizado:
$\frac12\langle 1s(1)2s(2) \pm 2s(1)1s(2)|\hat V_{12}|1s(1)2s(2) \pm
2s(1)1s(2)\rangle$. Los dos términos directos dan $\frac12(J + J)$; los dos términos
cruzados dan $\pm\frac12(K + K)$, ya que $\hat V_{12}$ es simétrico en los
electrones. Así, $\Delta E = J \pm K$, y las funciones de espín, normalizadas y
no afectadas por $\hat V_{12}$, solo multiplican por 1.
\end{proof}

El triplete queda más bajo aunque el hamiltoniano no contiene el espín: su
función espacial se anula cuando $\mathbf r_1 = \mathbf r_2$, de modo que los dos
electrones se evitan y se repelen menos. Este ``hueco de Fermi'' es el
contenido físico de la primera regla de Hund.

\begin{example}[La integral de intercambio del helio, medida]\label{ex:b3:many-electron-atoms:K}
% ledger: lev:He.2s3S1, lev:He.2s1S0
Los niveles $1s2s$ del helio están a \qty{159856}{cm^{-1}} (${}^3$S) y a
\qty{166277}{cm^{-1}} (${}^1$S) por encima del estado fundamental: el singlete está
\qty{6421}{cm^{-1}}, \qty{0.796}{eV}, más arriba, y $K(1s,2s) = \qty{0.398}{eV}$.
La figura siguiente muestra las dos familias.
\end{example}

\section{Acoplamiento de Russell--Saunders y símbolos de término}

En un átomo ligero, los momentos angulares orbitales de los electrones se suman en un
total $\mathbf L$, sus espines en un total $\mathbf S$, y solo después
$\mathbf L$ y $\mathbf S$ se acoplan débilmente entre sí.

\begin{definition}[Acoplamiento de Russell--Saunders]\label{def:b3:many-electron-atoms:russell-saunders}
En el \emph{acoplamiento de Russell--Saunders}\index{acoplamiento de Russell--Saunders}, los
estados de una configuración se etiquetan con el \emph{número cuántico de momento angular
orbital total}\index{numero cuantico de momento angular orbital total@número cuántico de momento angular orbital total}
$L$ ($\hat{\mathbf L}^2 = L(L+1)\hbar^2$, con $M_L = \sum m_l$), el \emph{número
cuántico de espín total}\index{numero cuantico de espin total@número cuántico de espín total} $S$ ($M_S = \sum m_s$)
y el \emph{número cuántico de momento angular total}\index{numero cuantico de momento angular total@número cuántico de momento angular total}
$J$, que toma los valores $|L - S|, \dots, L + S$.
\end{definition}

\begin{definition}[Término espectroscópico, multiplicidad de espín, símbolo de término]\label{def:b3:many-electron-atoms:term}
Un \emph{término espectroscópico}\index{termino espectroscopico@término espectroscópico} es el conjunto de los
$(2L+1)(2S+1)$ estados de una configuración con $L$ y $S$ dados. Su
\emph{multiplicidad de espín}\index{multiplicidad de espin@multiplicidad de espín} es $2S + 1$. El \emph{símbolo de término}
\index{simbolo de termino@símbolo de término} es ${}^{2S+1}L$, con las letras S, P, D, F, G, H
para $L = 0, 1, 2, 3, 4, 5$; un estado de $J$ dado se escribe ${}^{2S+1}L_J$
---por ejemplo, \termsym{3}{P}{2}.
\end{definition}

\begin{proposition}[Capas cerradas]\label{prop:b3:many-electron-atoms:closed-shell}
Una subcapa llena tiene $L = 0$ y $S = 0$; no contribuye en nada al
\omterm{def:b3:many-electron-atoms:term}{símbolo de término}.
\end{proposition}

\begin{proof}
En una subcapa llena, cada $m_l$ aparece dos veces y cada $m_s = \pm\frac12$
aparece $2l+1$ veces, luego $M_L = 0$ y $M_S = 0$; y solo hay una manera de
llenarla (un determinante). Un único estado con $M_L = M_S = 0$ solo puede
pertenecer a $L = 0$, $S = 0$.
\end{proof}

\begin{proposition}[Recuento de los estados de una configuración]\label{prop:b3:many-electron-atoms:count}
$N$ electrones en una subcapa de $2(2l+1)$ \omterm{def:b3:many-electron-atoms:spin-orbital}{espín-orbitales} dan
$\binom{2(2l+1)}{N}$ determinantes, y las \omterm{def:b3:quantum-model-systems:degenerate}{degeneraciones} $(2L+1)(2S+1)$ de
sus términos suman este número.
\end{proposition}

\begin{proof}
Un determinante queda fijado por el conjunto de los \omterm{def:b3:many-electron-atoms:spin-orbital}{espín-orbitales} ocupados (su orden
solo cambia su signo): se eligen $N$ de los $2(2l+1)$. Los términos son los
mismos estados reagrupados, de modo que los recuentos deben coincidir.
\end{proof}

\begin{method}[Términos de una configuración]\label{met:b3:many-electron-atoms:terms}
\begin{enumerate}
  \item Enumerar los determinantes permitidos por el principio de Pauli y tabularlos
        según $M_L$ y $M_S$.
  \item Buscar el mayor $M_L$; con el mayor $M_S$ que lo acompañe,
        inicia un término con $L = M_L$, $S = M_S$.
  \item Tachar un determinante por cada par $(M_L, M_S)$ con
        $|M_L| \le L$, $|M_S| \le S$.
  \item Repetir con lo que queda, hasta vaciar la tabla; comprobar el recuento.
\end{enumerate}
\end{method}

\begin{example}[Los términos de $p^2$]\label{ex:b3:many-electron-atoms:p2}
Dos electrones en $p$ ($m_l = 1, 0, -1$) dan $\binom62 = 15$ determinantes.
Su tabla según $M_L$ y $M_S$ es:
\begin{center}\small
\begin{tabular}{c|ccc}
\hline
 & $M_S = 1$ & $M_S = 0$ & $M_S = -1$\\ \hline
$M_L = 2$ & -- & $(1^+,1^-)$ & --\\
$M_L = 1$ & $(1^+,0^+)$ & $(1^+,0^-)$, $(1^-,0^+)$ & $(1^-,0^-)$\\
$M_L = 0$ & $(1^+,-1^+)$ & $(1^+,-1^-)$, $(1^-,-1^+)$, $(0^+,0^-)$ & $(1^-,-1^-)$\\
$M_L = -1$ & $(0^+,-1^+)$ & $(0^+,-1^-)$, $(0^-,-1^+)$ & $(0^-,-1^-)$\\
$M_L = -2$ & -- & $(-1^+,-1^-)$ & --\\ \hline
\end{tabular}
\end{center}
(con $m_l^\pm$ para $m_s = \pm\frac12$). $M_L = 2$ solo aparece con $M_S = 0$:
un término ${}^1$D (5 estados). El mayor valor restante, $M_L = 1$, va con
$M_S = 1$: un término ${}^3$P (9 estados). Queda un estado con $M_L = M_S = 0$: un
término ${}^1$S. En efecto, $5 + 9 + 1 = 15$.
\end{example}

\begin{example}[El carbono, medido]\label{ex:b3:many-electron-atoms:carbon}
% ledger: lev:C.3P1, lev:C.3P2, lev:C.1D2, lev:C.1S0
La configuración fundamental $2p^2$ del carbono da los niveles \termsym{3}{P}{0},
\termsym{3}{P}{1}, \termsym{3}{P}{2} a 0, 16.4 y \qty{43.4}{cm^{-1}}, luego
\termsym{1}{D}{2} a \qty{10193}{cm^{-1}} y \termsym{1}{S}{0} a
\qty{21648}{cm^{-1}}: el término ${}^3$P es el más bajo, como predicen las reglas de Hund.
\end{example}

\begin{omfigure}
\begin{tikzpicture}
\begin{axis}[name=main, width=0.62\linewidth, height=6.4cm, xmin=-0.3, xmax=6.6,
  ymin=-1500, ymax=23500, axis y line=left, axis x line=none,
  ylabel={energía / cm$^{-1}$}, unbounded coords=jump, clip=false,
  scaled y ticks=false, ytick={0,10000,20000},
  tick label style={font=\small}, label style={font=\small}]
\addplot[omstate] table[x=x, y=E] {figdata/out/bachelor-3-atomic-levels-carbon.dat};
\draw[densely dotted, gray] (axis cs:1,4858.5) -- (axis cs:2,29.6);
\draw[densely dotted, gray] (axis cs:1,4858.5) -- (axis cs:2,10192.7);
\draw[densely dotted, gray] (axis cs:1,4858.5) -- (axis cs:2,21648);
\draw[densely dotted, gray] (axis cs:3,10192.7) -- (axis cs:4,10192.7);
\draw[densely dotted, gray] (axis cs:3,21648) -- (axis cs:4,21648);
\node[font=\small, anchor=south] at (axis cs:0.5,4858.5) {$2p^2$};
\node[font=\small, anchor=south] at (axis cs:2.5,21648) {${}^1$S};
\node[font=\small, anchor=south] at (axis cs:2.5,10192.7) {${}^1$D};
\node[font=\small, anchor=south] at (axis cs:2.5,29.6) {${}^3$P};
\node[font=\small, anchor=west] at (axis cs:5.05,21648) {\termsym{1}{S}{0}};
\node[font=\small, anchor=west] at (axis cs:5.05,10192.7) {\termsym{1}{D}{2}};
\node[font=\scriptsize, anchor=north] at (axis cs:0.5,-600) {configuración};
\node[font=\scriptsize, anchor=north] at (axis cs:2.5,-600) {términos};
\node[font=\scriptsize, anchor=north] at (axis cs:4.5,-600) {niveles};
\draw[gray, ->] (axis cs:3.05,60) -- (axis cs:6.9,2500);
\end{axis}
\begin{axis}[at={(main.south east)}, anchor=south west, xshift=1.2cm, yshift=0.9cm,
  width=3.6cm, height=3.6cm, xmin=-0.1, xmax=1.1, ymin=-5, ymax=50,
  axis y line=right, axis x line=none, unbounded coords=jump, ytick={0,16.4,43.4},
  yticklabels={0,16.4,43.4}, tick label style={font=\scriptsize}, clip=false,
  title={\scriptsize ${}^3$P, ampliado}, title style={yshift=-4pt}]
\addplot[omstate] table[x=x, y=E] {figdata/out/bachelor-3-atomic-levels-carbon-zoom.dat};
\node[font=\scriptsize, anchor=east] at (axis cs:0,0) {$J = 0$};
\node[font=\scriptsize, anchor=east] at (axis cs:0,16.4) {$J = 1$};
\node[font=\scriptsize, anchor=east] at (axis cs:0,43.4) {$J = 2$};
\end{axis}
\end{tikzpicture}

{\small La configuración fundamental del carbono: una configuración, tres términos
(en sus energías medidas; el término ${}^3$P en la media ponderada de sus
niveles, la configuración en la media de los quince estados) y cinco
niveles. El desdoblamiento espín--órbita de ${}^3$P (ampliado, en
\unit{cm^{-1}}) es unos cientos de veces menor que las separaciones entre
términos.}
\end{omfigure}

\section{Reglas de Hund y acoplamiento espín--órbita}

\begin{proposition}[Reglas de Hund para los términos]\label{prop:b3:many-electron-atoms:hund-terms}
Para la configuración fundamental de un átomo o ion, el término más bajo es el que tiene
(1) el mayor $S$ y, después, (2) para ese $S$, el mayor $L$. Su nivel más bajo
tiene (3) $J = |L - S|$ si la subcapa está menos que semillena, y $J = L +
S$ si está más que semillena.
\end{proposition}
\begin{proof}[Estatus]
Estas reglas están establecidas experimentalmente, solo para las configuraciones fundamentales;
no son teoremas. La regla 1 es la estabilización por intercambio de la sección
anterior; la regla 2 dice que los electrones que giran en el mismo sentido se encuentran
con menos frecuencia; la regla 3 se deduce del signo del \omterm{def:b3:many-electron-atoms:spin-orbit}{acoplamiento espín--órbita}, más abajo.
\end{proof}

\begin{method}[El término fundamental a partir de un diagrama de casillas]\label{met:b3:many-electron-atoms:ground-term}
\begin{enumerate}
  \item Llenar las casillas de la subcapa abierta desde $m_l = +l$ hacia abajo, un
        electrón por casilla con espines paralelos, y aparear después de nuevo desde $+l$.
  \item $S = \frac12 \times$(número de electrones desapareados); $L = |\sum m_l|$.
  \item $J = |L - S|$ si está menos que semillena, $L + S$ si está más que
        semillena; si está exactamente semillena, $L = 0$ y $J = S$.
\end{enumerate}
\end{method}

\begin{example}[Términos fundamentales]\label{ex:b3:many-electron-atoms:ground-terms}
% ledger: lev:Ti2.3P0, lev:Ni2.3P0, lev:Cr3.4P1_2
\ce{Ti^2+}, $d^2$: casillas $m_l = 2, 1$ llenas, $S = 1$, $L = 3$, $J = 2$:
\termsym{3}{F}{2}. \ce{Cr^3+}, $d^3$: $S = \frac32$, $L = 2 + 1 + 0 = 3$,
\termsym{4}{F}{3/2}. \ce{Fe^2+}, $d^6$: cinco hacia arriba y uno hacia abajo en $m_l = 2$:
$S = 2$, $L = 2$, más que semillena, \termsym{5}{D}{4}. \ce{Ni^2+}, $d^8$:
$S = 1$, $L = 3$, \termsym{3}{F}{4}. Todos coinciden con el nivel fundamental que
mide la espectroscopia atómica.
\end{example}

\begin{definition}[Acoplamiento espín--órbita, estructura fina]\label{def:b3:many-electron-atoms:spin-orbit}
El \emph{acoplamiento espín--órbita}\index{acoplamiento espin--orbita@acoplamiento espín--órbita} es la interacción
entre los momentos magnéticos del espín de un electrón y de su movimiento
orbital, que se escribe $A\,\hat{\mathbf L}\cdot\hat{\mathbf S}/\hbar^2$ para un término.
Desdobla un término en sus niveles de distinto $J$: la \emph{estructura
fina}\index{estructura fina} del término.
\end{definition}

\begin{proposition}[Regla de los intervalos de Landé]\label{prop:b3:many-electron-atoms:lande}
Dentro de un término, $E(J) = E_0 + \frac{A}{2}[J(J+1) - L(L+1) - S(S+1)]$, de modo que
$E(J) - E(J - 1) = AJ$. $A > 0$ (normal) para una subcapa menos que semillena,
$A < 0$ (invertida) para una más que semillena.
\end{proposition}

\begin{proof}
$\hat{\mathbf J} = \hat{\mathbf L} + \hat{\mathbf S}$ da $\hat{\mathbf J}^2 =
\hat{\mathbf L}^2 + \hat{\mathbf S}^2 + 2\hat{\mathbf L}\cdot\hat{\mathbf S}$, luego
$\hat{\mathbf L}\cdot\hat{\mathbf S} = \frac12(\hat{\mathbf J}^2 - \hat{\mathbf L}^2 -
\hat{\mathbf S}^2)$, cuyo \omterm{def:b3:quantum-model-systems:operator}{valor propio} en un estado de $J$, $L$, $S$ dados es
$\frac{\hbar^2}{2}[J(J+1) - L(L+1) - S(S+1)]$. La diferencia entre $J$ y
$J - 1$ es $\frac A2[J(J+1) - (J-1)J] = AJ$. El signo de $A$ se admite: una
subcapa más que semillena se comporta como el número correspondiente de huecos
positivos.
\end{proof}

\begin{example}[Puesta a prueba de la regla de los intervalos]\label{ex:b3:many-electron-atoms:lande-test}
% ledger: lev:C.3P1, lev:C.3P2, lev:O.3P1, lev:O.3P0
Para ${}^3$P la regla predice intervalos en la razón $J = 2 : 1$, es decir, 2.
Carbono: $(43.4 - 16.4)/16.4 = 1.65$. Oxígeno ($2p^4$, invertido, \termsym{3}{P}{2}
el más bajo, luego $J = 1$ a \qty{158.3}{cm^{-1}} y $J = 0$ a \qty{227.0}{cm^{-1}}):
$158.3/68.7 = 2.30$. La regla se cumple con un margen del 20\,\%: el \omterm{def:b3:many-electron-atoms:russell-saunders}{acoplamiento de Russell--Saunders}
es una buena descripción de los átomos ligeros, no una descripción exacta.
\end{example}

\begin{omfigure}
\begin{tikzpicture}[font=\small]
  \draw[omstate] (0,0) -- (1.6,0) node[midway, below] {$3s\ {}^2$S$_{1/2}$};
  \draw[omstate] (0,3.0) -- (1.6,3.0);
  \node[anchor=east] at (-0.1,3.0) {$3p$};
  \draw[omstate] (2.8,2.75) -- (4.4,2.75) node[right] {${}^2$P$_{1/2}$};
  \draw[omstate] (2.8,3.25) -- (4.4,3.25) node[right] {${}^2$P$_{3/2}$};
  \draw[densely dotted, gray] (1.6,3.0) -- (2.8,2.75) (1.6,3.0) -- (2.8,3.25);
  \draw[omfluo] (3.3,2.73) -- (3.3,0.02) node[pos=0.55, left] {D$_1$};
  \draw[omfluo] (3.9,3.23) -- (3.9,0.02) node[pos=0.55, right] {D$_2$};
  \draw[omstate] (2.8,0) -- (4.4,0);
  \draw[densely dotted, gray] (1.6,0) -- (2.8,0);
  \draw[decorate, decoration={brace, amplitude=3pt}] (5.5,3.27) -- (5.5,2.73)
     node[midway, right=3pt, align=left] {\footnotesize \qty{17.2}{cm^{-1}}\\[-1pt]\footnotesize (no a escala)};
\end{tikzpicture}

{\small Las rayas D del sodio. El nivel $3p$ se desdobla por \omterm{def:b3:many-electron-atoms:spin-orbit}{acoplamiento espín--órbita}
en ${}^2$P$_{1/2}$ y ${}^2$P$_{3/2}$; la emisión hacia el nivel fundamental $3s$ da
dos rayas amarillas, D$_1$ a \qty{589.76}{nm} y D$_2$ a \qty{589.16}{nm} (longitudes
de onda en el vacío).}
\end{omfigure}

\begin{example}[El doblete del sodio]\label{ex:b3:many-electron-atoms:sodium}
% ledger: lev:Na.3p1_2, lev:Na.3p3_2
Los niveles $3p$ del sodio están a 16956.17 y \qty{16973.37}{cm^{-1}}: las rayas
D están a $10^7/16956.17 = \qty{589.76}{nm}$ y \qty{589.16}{nm} en el vacío,
separadas por \qty{17.20}{cm^{-1}}. Para ${}^2$P, $E(\frac32) - E(\frac12) =
\frac32A$, luego $A = \qty{11.47}{cm^{-1}}$.
\end{example}

\begin{proposition}[Reglas de selección para los átomos]\label{prop:b3:many-electron-atoms:selection-atoms}
En el \omterm{def:b3:many-electron-atoms:russell-saunders}{acoplamiento de Russell--Saunders}, una transición dipolar eléctrica exige
$\Delta S = 0$, $\Delta L = 0, \pm1$, $\Delta J = 0, \pm1$ (pero no $J = 0 \to J
= 0$) y un cambio de paridad: para el salto de un solo electrón, $\Delta l = \pm1$.
\end{proposition}
\begin{proof}[Estatus]
$\Delta S = 0$ se demuestra en el \cref{ch:b3:electronic-spectroscopy}: el \omterm{def:b3:quantum-model-systems:operator}{operador}
dipolar no actúa sobre el espín. Las demás se deducen del carácter vectorial
del \omterm{def:b3:quantum-model-systems:operator}{operador} dipolar y aquí se admiten.
\end{proof}

\begin{omfigure}
\begin{tikzpicture}
\begin{axis}[width=0.62\linewidth, height=6.2cm, xmin=-0.5, xmax=3.5,
  ymin=158, ymax=173, axis y line=left, axis x line=none,
  ylabel={energía / $10^3$ cm$^{-1}$}, unbounded coords=jump, clip=false,
  tick label style={font=\small}, label style={font=\small}]
\addplot[omstate] table[x=x, y=E] {figdata/out/bachelor-3-atomic-levels-helium.dat};
\node[font=\small, anchor=west] at (axis cs:1.05,166.277) {$1s2s\ {}^1$S};
\node[font=\small, anchor=west] at (axis cs:1.05,171.135) {$1s2p\ {}^1$P};
\node[font=\small, anchor=west] at (axis cs:3.05,159.856) {$1s2s\ {}^3$S};
\node[font=\small, anchor=west] at (axis cs:3.05,169.087) {$1s2p\ {}^3$P};
\node[font=\small, anchor=south] at (axis cs:0.5,172.3) {singletes};
\node[font=\small, anchor=south] at (axis cs:2.5,172.3) {tripletes};
\draw[omfluo] (axis cs:0.5,171.0) -- (axis cs:0.5,166.42) node[pos=0.5, left, font=\scriptsize] {\qty{2059}{nm}};
\draw[omfluo] (axis cs:2.5,168.95) -- (axis cs:2.5,160.0) node[pos=0.5, left, font=\scriptsize] {\qty{1083}{nm}};
\draw[densely dotted, gray] (axis cs:-0.3,159.856) -- (axis cs:2,159.856);
\draw[<->, omThm] (axis cs:-0.2,159.856) -- (axis cs:-0.2,166.277) node[pos=0.5, right, font=\scriptsize] {$2K$};
\end{axis}
\end{tikzpicture}

{\small Los dos ``helios''. Niveles singlete y triplete de las configuraciones
$1s2s$ y $1s2p$ en sus energías medidas por encima del estado fundamental (el
término ${}^3$P en la media de sus niveles). Las líneas solo unen estados de la misma
familia; cada triplete queda por debajo de su singlete, en $2K$.}
\end{omfigure}

\begin{history}[Dos helios y el principio de exclusión]
El helio recibió su nombre del Sol en 1868 y William Ramsay lo encontró en la Tierra
en 1895. Sus dos series de rayas desconcertaron a los espectroscopistas durante treinta
años. En 1925 Wolfgang Pauli propuso que dos electrones no pueden compartir los
mismos números cuánticos, el mismo año en que George Uhlenbeck y Samuel Goudsmit
propusieron el espín del electrón; en 1926 Werner Heisenberg demostró que la
simetría de la función de onda basta por sí sola para separar el helio en sus familias singlete y
triplete.
\end{history}

\begin{inthelab}[Resolver el doblete del sodio]
Se coloca una lámpara de sodio ante la rendija de un espectrómetro de red. Con una
red de 600 líneas por milímetro, los ángulos de primer orden de las dos rayas
difieren en unos \num{0.02}\textdegree: un buen espectrómetro las separa, un
prisma de prácticas no. La lámpara funciona muy caliente; solo se manipula cuando se ha
enfriado.
\end{inthelab}

\section{Ejercicios}

\begin{exercise}[$\star$]\label{exo:b3:many-electron-atoms:1}
Escríbase el \omterm{def:b3:many-electron-atoms:slater}{determinante de Slater} del estado fundamental del litio, $1s^22s$, con
el electrón $2s$ de espín $\alpha$. ¿Por qué no existe determinante para
$1s^3$?
\end{exercise}

\begin{exercise}[$\star$]\label{exo:b3:many-electron-atoms:2}
Cuéntense los determinantes de las configuraciones $p^2$, $p^3$ y $d^2$. Los términos
de $p^3$ son ${}^4$S, ${}^2$D, ${}^2$P y los de $d^2$ son ${}^3$F, ${}^3$P, ${}^1$G,
${}^1$D, ${}^1$S: compruébense ambos recuentos.
\end{exercise}

\begin{exercise}[$\star$]\label{exo:b3:many-electron-atoms:3}
Dense el término y el nivel fundamentales de N, O, F, \ce{Ti^2+} y \ce{Fe^3+} ($d^5$).
\end{exercise}

\begin{exercise}[$\star$]\label{exo:b3:many-electron-atoms:4}
Enumérense los niveles de los términos ${}^2$D y ${}^3$P con sus \omterm{def:b3:quantum-model-systems:degenerate}{degeneraciones} $2J+1$,
y compruébese que las \omterm{def:b3:quantum-model-systems:degenerate}{degeneraciones} suman $(2L+1)(2S+1)$.
\end{exercise}

\begin{exercise}[$\star\star$]\label{exo:b3:many-electron-atoms:5}
Los niveles ${}^3$P del carbono están a 0, 16.4 y \qty{43.4}{cm^{-1}}. Calcúlense la
razón de los intervalos y la constante espín--órbita $A$ a partir de cada intervalo.
¿Qué indica la diferencia entre los dos valores de $A$?
\end{exercise}

\begin{exercise}[$\star\star$]\label{exo:b3:many-electron-atoms:6}
Los niveles $3p$ del sodio están a 16956.17 y \qty{16973.37}{cm^{-1}}. Calcúlense
las longitudes de onda en el vacío de las rayas D, su separación en \unit{cm^{-1}} y
en meV, y $A$ para el término $3p$.
\end{exercise}

\begin{exercise}[$\star\star$]\label{exo:b3:many-electron-atoms:7}
¿Cuáles de estas transiciones del sodio están permitidas en emisión: $3p \to 3s$,
$3d \to 3s$, $3d \to 3p$, $4s \to 3s$, $4s \to 3p$? Justifíquese cada una.
\end{exercise}

\begin{exercise}[$\star\star$]\label{exo:b3:many-electron-atoms:8}
% ledger: lev:He.2p3P2, lev:He.2p3P1, lev:He.2p3P0, lev:He.2p1P1
Los niveles $1s2p$ del helio son \termsym{3}{P}{2}, \termsym{3}{P}{1},
\termsym{3}{P}{0} a 169086.76, 169086.84, \qty{169087.83}{cm^{-1}} y
\termsym{1}{P}{1} a \qty{171134.89}{cm^{-1}}. Calcúlese la energía media del
término ${}^3$P y después $K(1s,2p)$ en eV. Compárese con $K(1s,2s) = \qty{0.398}{eV}$
y explíquese la diferencia.
\end{exercise}

\begin{exercise}[$\star\star$]\label{exo:b3:many-electron-atoms:9}
Encuéntrense los términos de $d^2$ por el método de este capítulo, empezando por el
mayor $M_L$ (puede omitirse escribir la tabla completa, pero justifíquese cada término).
\end{exercise}

\begin{exercise}[$\star\star\star$]\label{exo:b3:many-electron-atoms:10}
Demuéstrese que la función espacial triplete $\frac{1}{\sqrt2}[1s(1)2s(2) -
2s(1)1s(2)]$ se anula cuando $\mathbf r_1 = \mathbf r_2$ y que la función singlete
no. Relaciónese con el signo de $E_+ - E_-$.
\end{exercise}

\begin{exercise}[$\star\star\star$]\label{exo:b3:many-electron-atoms:11}
% ledger: lev:Ni2.3F3, lev:Ni2.3F2
Para \ce{Ni^2+} ($d^8$), los niveles ${}^3$F están a 0 ($J = 4$), 1360.7 ($J = 3$)
y \qty{2269.6}{cm^{-1}} ($J = 2$). Explíquese el orden, calcúlese la razón de
los intervalos y compárese con la predicción de Landé. Dese $A$ a partir del intervalo
mayor.
\end{exercise}

\begin{exercise}[$\star\star\star$]\label{exo:b3:many-electron-atoms:12}
Demuéstrese que el \omterm{def:b3:quantum-model-systems:hermitian}{valor esperado} de $\hat{\mathbf L}\cdot\hat{\mathbf S}$, sumado
sobre todos los $(2L+1)(2S+1)$ estados de un término, es nulo, de modo que el \omterm{def:b3:many-electron-atoms:spin-orbit}{acoplamiento
espín--órbita} deja inalterada la media ponderada de los niveles. Compruébese en el
${}^3$P del carbono.
\end{exercise}

\section{Problema: dos helios}

\begin{problem}[{Problema de fin de semana --- las familias singlete y triplete del
helio: determinantes, los términos de los iones de metales de transición, las reglas
de selección que separan las dos familias y la integral de intercambio medida}]
\label{pb:b3:many-electron-atoms:1}
% ledger: lev:He.2s3S1, lev:He.2s1S0, lev:He.2p3P2, lev:He.2p3P1, lev:He.2p3P0, lev:He.2p1P1, lev:Ti2.3F3, lev:Ti2.3F4, lev:Ti2.3P0, lev:Ti2.3P1, lev:Ti2.3P2
Datos (NIST, en \unit{cm^{-1}} por encima del estado fundamental de cada especie). Helio:
$1s2s$ \termsym{3}{S}{1} 159855.97, \termsym{1}{S}{0} 166277.44; $1s2p$
\termsym{3}{P}{2,1,0} 169086.76, 169086.84, 169087.83, \termsym{1}{P}{1}
171134.89. \ce{Ti^2+}: \termsym{3}{F}{2,3,4} 0, 184.9, 420.4; \termsym{3}{P}{0,1,2}
10538.4, 10603.6, 10721.2. $\qty{1}{eV} = \qty{8065.54}{cm^{-1}}$.

\medskip\noindent\textbf{Parte I --- \omterm{def:b3:many-electron-atoms:spin-orbital}{Espín-orbitales} y determinantes.}
\begin{enumerate}
  \item Enumérense los cuatro \omterm{def:b3:many-electron-atoms:spin-orbital}{espín-orbitales} disponibles para la configuración $1s2s$.
  \item Escríbanse los cuatro determinantes con un electrón en $1s$ y otro en
        $2s$.
  \item Demuéstrese que $|1s\alpha\,2s\alpha|$ es el producto de una función espacial
        antisimétrica por la función de espín $\alpha(1)\alpha(2)$.
  \item Combínense $|1s\alpha\,2s\beta|$ y $|1s\beta\,2s\alpha|$ para obtener la
        componente $M_S = 0$ de la misma función espacial, y un cuarto estado.
  \item ¿Qué función espacial, simétrica o antisimétrica, acompaña a
        $S = 0$? ¿Y a $S = 1$?
  \item ¿Por qué se llaman singlete y triplete?
\end{enumerate}

\medskip\noindent\textbf{Parte II --- Los términos de \ce{Ti^2+}.}
\begin{enumerate}[resume]
  \item Dense la configuración de \ce{Ti^2+} y el número de sus
        determinantes.
  \item Encuéntrense su término y su nivel fundamentales con las reglas de Hund; compruébense con los
        datos.
  \item Sus otros términos son ${}^3$P, ${}^1$G, ${}^1$D y ${}^1$S: compruébese el recuento de
        estados.
  \item ¿Es normal o invertida la \omterm{def:b3:many-electron-atoms:spin-orbit}{estructura fina} de ${}^3$F? Póngase a prueba la regla de los
        intervalos de Landé.
  \item Calcúlense las energías medias ponderadas de los términos ${}^3$F y ${}^3$P.
  \item Su separación es $15B$, donde $B$ es un parámetro de Racah del
        ion (\cref{ch:b3:complex-spectra-magnetism}). Calcúlese $B$.
\end{enumerate}

\medskip\noindent\textbf{Parte III --- Por qué dos familias.}
\begin{enumerate}[resume]
  \item ¿Qué regla de selección prohíbe una transición entre un singlete y un
        triplete?
  \item El triplete más bajo, $1s2s$ \termsym{3}{S}{1}, no puede emitir hacia el estado
        fundamental $1s^2$ \termsym{1}{S}{0}. Dense dos razones. ¿Cómo se llama un estado
        así?
  \item Calcúlese la longitud de onda de la raya $1s2p\ {}^3$P $\to 1s2s\ {}^3$S (úsese
        la media de los niveles ${}^3$P).
  \item Calcúlese la longitud de onda de $1s2p\ {}^1$P $\to 1s2s\ {}^1$S.
  \item Calcúlese la longitud de onda de $1s2p\ {}^1$P $\to 1s^2\ {}^1$S. ¿En qué
        región del espectro está?
  \item Explíquese por qué los espectroscopistas del siglo XIX, al ver estas rayas,
        creyeron en dos gases distintos.
\end{enumerate}

\medskip\noindent\textbf{Parte IV --- La \omterm{def:b3:many-electron-atoms:exchange}{integral de intercambio}, medida.}
\begin{enumerate}[resume]
  \item Escríbanse las energías de primer orden del singlete y del triplete $1s2s$ en
        función de $E_{1s}$, $E_{2s}$, $J$ y $K$.
  \item ¿Cuál es más bajo y por qué, en términos de las posiciones de los electrones?
  \item Calcúlese la separación singlete--triplete $1s2s$ en \unit{cm^{-1}} y en eV.
  \item Calcúlese $K(1s,2p)$ en eV a partir de los niveles $1s2p$.
  \item ¿Por qué $K(1s,2s)$ es mayor que $K(1s,2p)$?
  \item ¿Confirman los estados excitados del helio la primera regla de Hund?
  \item Enúnciese el resultado: la \omterm{def:b3:many-electron-atoms:exchange}{integral de intercambio} $K(1s,2s)$ del helio, en
        eV.
\end{enumerate}
\end{problem}
